\chapter{Discusión y conclusiones}
\label{cap:discusion}

\section{Discusión de los resultados}

\subsection{Seguimiento de precios y capacidad de arbitraje (OE2, H1)}

El resultado de H1 es el más sólido del trabajo: el arbitraje activo reduce el área de exceso de error en todas las semillas y en todas las celdas, y la capacidad gradúa la recuperación, de $4.6\,\%$ sin arbitraje a $56.8\,\%$ con capacidad baja y a $100\,\%$ con capacidad alta. Es un resultado esperable desde la mecánica del modelo, porque el arbitrajista es el único agente que mueve el precio del AMM hacia la referencia, y por eso la contribución no está en el signo del efecto sino en su graduación. La tasa de recuperación con capacidad baja depende de la combinación entre el tamaño del choque y la profundidad (tabla~\ref{tab:recuperacion}): un choque grande requiere varios pasos acotados, y el horizonte de 20 ticks puede no bastar. Esa dependencia sugiere que, en este modelo, la capacidad de arbitraje relevante no es un umbral sino una razón entre el tamaño de la discrepancia y la capacidad por tick.

La lección de diseño es la precisión. La semiamplitud objetivo de H1 no se alcanzó porque el piloto combinó choques positivos y negativos cuyos efectos se cancelan casi por completo dentro de cada semilla (correlación $-0.981$). Un piloto construido con el mismo diseño que la confirmación habría anticipado una desviación nueve veces mayor. La conclusión de que el arbitraje reduce el exceso no depende de ello, pero la estimación del tamaño del efecto es menos precisa de lo previsto.

\subsection{Liquidez adaptativa (OE3, H2)}

H2 se sostiene en promedio, pero no de manera uniforme. El retiro se activó en todas las ejecuciones de E3 y, sin embargo, la respuesta máxima solo aumenta con la menor profundidad. Con las profundidades 1 y 2 el máximo es el mismo en E1-alta y en E3. Con la profundidad 0.5, la comprobación exploratoria sin agente de pánico conserva aproximadamente la mitad del efecto, de modo que este depende del retiro y, en igual medida, de la interacción entre el retiro y el agente de pánico. Lo que el experimento permite afirmar es que, en este micromundo, una regla de retiro de un solo disparo no amplifica la respuesta máxima cuando las órdenes son pequeñas respecto de la reserva, y que la amplificación aparece cuando la piscina es suficientemente superficial para que una orden de ruido active la reacción de pánico.

No se afirma que los proveedores reales amplifiquen las perturbaciones. La regla de retiro es una decisión de modelado, sin calibración y sin respaldo empírico propio. Su respaldo teórico es indirecto: la literatura trata la salida de los proveedores como un problema de decisión \cite{bergault2025exit} y documenta su comportamiento en intercambios descentralizados \cite{lehar2025uniswap}, pero ninguno de esos trabajos valida la regla usada aquí.

\subsection{Mapa dinámico reducido (OE4, H3)}

El resultado de H3 es mixto. El mapa de dos variables predice mejor que la persistencia, y con una mejora sustancial en las condiciones con arbitraje activo (15 a 21 por ciento en E1-alta y E3), pero no supera a un autorregresivo escalar y es ligeramente peor que la persistencia sin arbitraje. Dos observaciones limitan su interpretación. Primera, el segundo estado, la presión neta de compra, no mejora la predicción a un paso respecto del exceso de error solo; la ventaja del mapa sobre la persistencia proviene del coeficiente de contracción $a$, que un modelo escalar captura igual. Segunda, el coeficiente $\beta$ de E3 tiene signo opuesto al de las demás condiciones y la ecuación del exceso es degenerada en E1-máx, lo que muestra que la estructura de dos variables no es estable entre condiciones.

La clasificación de regímenes con el criterio de Schur no discrimina: todas las matrices ajustadas son estables, incluida la de E0, cuyo radio espectral es $0.998$ y cuyas ejecuciones son persistentes en $57\,\%$ de los casos. Esto es coherente con lo que el criterio puede decir (la estabilidad de una dinámica lineal autónoma en un punto) y no con el horizonte finito y el ruido del simulador. El mapa sirve como descripción de la tasa de contracción del error, no como clasificador de regímenes.

\section{Conclusiones por objetivo}

\begin{description}
  \item[OE1.] Se formuló y verificó un simulador discreto con nueve invariantes y un registro reproducible: la ejecución confirmatoria se regenera con identificador y sumas de verificación idénticos, salvo la marca de tiempo. La verificación demuestra que la implementación respeta su especificación, no que la especificación sea realista.
  \item[OE2.] El arbitraje activo reduce el área de exceso de error de $2.36$ a $0.10$, y la capacidad gradúa la tasa de recuperación en las condiciones simuladas.
  \item[OE3.] El retiro adaptativo aumenta la respuesta máxima en $0.011$ en promedio, pero solo con la menor profundidad ensayada, y aproximadamente la mitad de ese efecto se atribuye al agente de pánico. No se concluye amplificación con profundidades mayores.
  \item[OE4.] El mapa reducido supera a la persistencia, no supera a un autorregresivo escalar, y su clasificación de regímenes fuera de muestra es débil. Es una descripción local, y su utilidad queda limitada a la tasa de contracción del error.
\end{description}

\section{Conclusión general}

En el micromundo estudiado, la capacidad de arbitraje es el determinante principal del seguimiento del precio de referencia por un AMM de producto constante. La profundidad modifica el impacto de las órdenes de forma analítica y, junto con una regla de reacción de los agentes, puede amplificar la respuesta. Un mapa de dos variables describe la contracción del error mejor que la persistencia, pero un modelo de una variable lo describe igual de bien. Estas conclusiones valen para el modelo y no se extienden a mercados reales.

\section{Limitaciones y amenazas a la validez}

\begin{itemize}
  \item \textbf{Referencia exógena.} El CEX no responde al AMM, ni tiene profundidad, latencia o comisiones. Los efectos de realimentación entre mercados quedan fuera.
  \item \textbf{Sin comisiones.} El simulador solo implementa intercambios sin comisión. La banda de no arbitraje es un parámetro y no una consecuencia de un costo. Los niveles de comisión, de latencia y de volatilidad externa son candidatos para trabajos posteriores.
  \item \textbf{Calendario.} El orden de las acciones dentro de un tick afecta los resultados (por ejemplo, el agente de pánico observa el precio antes de que actúen el proveedor y el arbitrajista). No se evaluaron calendarios alternativos.
  \item \textbf{Definición de profundidad.} La profundidad escala las reservas y mantiene constante el tamaño absoluto de las órdenes. Otra definición, con órdenes proporcionales a la reserva, podría cambiar la conclusión sobre H2.
  \item \textbf{Inventarios.} Los inventarios son amplios y ninguna orden estuvo limitada por capital; las restricciones de capital reales no se representan.
  \item \textbf{Agentes de un solo tipo.} Cada clase de agente es representativa y no hay heterogeneidad dentro de una clase.
  \item \textbf{Solo datos sintéticos.} No hay datos de mercado ni validación externa. La coherencia entre las series simuladas y los hechos estilizados de mercados reales no se evaluó.
  \item \textbf{Comprobación exploratoria.} La repetición de H2 sin agente de pánico no estaba preregistrada.
  \item \textbf{Generalización de H3.} La evaluación fuera de muestra cubre celdas y semillas no vistas, pero dentro del mismo simulador y de la misma familia de choques.
\end{itemize}
